\documentclass[11pt,a4paper]{article}
\usepackage[spanish]{babel}
\usepackage[utf8]{inputenc}
\usepackage[T1]{fontenc}
\usepackage{amsmath,amssymb}
\usepackage[margin=2.7cm]{geometry}
\title{Taller de \LaTeX\\Algoritmos y Estructuras de Datos I}
\author{Fernando Toralez}
\date{}
\newcommand{\ent}{\mathbb Z}
\newcommand{\seq}[1]{\operatorname{seq}\langle#1\rangle}
\newcommand{\longitud}[1]{\lvert#1\rvert}
\begin{document}
\maketitle
\tableofcontents
\newpage
\section{Productos y factorial}
El factorial de un entero positivo se define como
\[
 n!=\prod_{i=1}^{n}i.
\]
Por ejemplo:
\[
 5!=\prod_{i=1}^{5}i=1\cdot2\cdot3\cdot4\cdot5=120.
\]
\section{Especificación}
\subsection{Factorial}
\begin{description}
 \item[Procedimiento] $\operatorname{factorial}(\mathrm{in}\ n:\ent,
              \mathrm{out}\ resultado:\ent)$.
 \item[Precondición] $n\geq0$.
 \item[Postcondición]
 $(n=0\Rightarrow resultado=1)\land
  (n>0\Rightarrow resultado=\prod_{k=1}^{n}k)$.
\end{description}
La especificación establece qué resultado se espera, sin fijar un algoritmo.
\subsection{Predicados sobre secuencias}
Definimos primero la primalidad para poder utilizarla en los predicados:
\[
 \operatorname{esPrimo}(n)\equiv n>1\land
 \neg\exists d\in\ent:\ 1<d<n\land n\bmod d=0.
\]
Para $s:\seq{\ent}$:
\begin{align*}
 \operatorname{todosPrimos}(s)
 &\equiv \forall i\in\ent:\ 0\leq i<\longitud{s}
                   \Rightarrow \operatorname{esPrimo}(s[i]),\\
 \operatorname{alMenosUnPrimo}(s)
 &\equiv \exists i\in\ent:\ 0\leq i<\longitud{s}
                   \land \operatorname{esPrimo}(s[i]).
\end{align*}
La secuencia vacía satisface el primero y no satisface el segundo.
\subsection{Suma de primos}
\[
 \operatorname{sumaPrimos}(s)=
 \sum_{i=0}^{\longitud{s}-1}
 \begin{cases}
   s[i],&\text{si }\operatorname{esPrimo}(s[i]),\\
   0,&\text{en otro caso}.
 \end{cases}
\]
La suma vacía se define como cero.
\section{Notación de un promedio cuadrático}
Como práctica adicional de acentos, índices y sumatorias:
\[
 \widehat R(\widehat f,\overline D)=
 \frac{1}{\lvert\overline D\rvert}
 \sum_{i:(X_i,Y_i)\in\overline D}
       (Y_i-\widehat f(X_i))^2,\qquad \lvert\overline D\rvert>0.
\]
Es un ejemplo de notación matemática; la no vacuidad evita dividir por cero.
\end{document}
